%% 2012-11-19改定（中鉢）
%% - 書式設定を「kiyou.sty」ファイルに移動しました

\documentclass[a4j,twoside]{jarticle}
\usepackage{kiyou} % AIIT紀要スタイルの読み込み
%% 必要なパッケージを追加してください
\usepackage{graphicx}
%\usepackage{latexsym}
%\usepackage{amsmath}
%\usepackage{amsthm}

%% 英語でキャプションを記載する場合は下記2行のコメントを外してください
% \renewcommand{\figurename}{Fig.}
% \renewcommand{\tablename}{Table.}

\begin{document}
\pagestyle{empty}
%%%%%%%%%%%%%%%%%%%%%%%%%%%%%%%%%%%%%%%%%%%%%%%%%%%%%%%%
%% 和文標題・英文標題・アブストラクト・キーワード
%%%%%%%%%%%%%%%%%%%%%%%%%%%%%%%%%%%%%%%%%%%%%%%%%%%%%%%%
\begin{center}
{\LARGE
%% 論文標題
産技大紀要フォーマットについて(2012-11-22版)
}\vskip\baselineskip{\large 
%% 和文著者名
産　技　太　郎$^{*}$・産　技　花　子$^{**}$
}\vskip\baselineskip{\Large 
%% 論文英語タイトル
Style and Layout of an AIIT Bulletin
}\vskip\baselineskip{\large 
%% 英文著者名
Taro Sangi$^{*}$ and Hanako Sangi$^{**}$
}\vskip\baselineskip
{\large Abstract}
\end{center}

{\normalsize\noindent
%% アブストラクト
Abstract in English. (approx. 100 Words) 
}\vskip\baselineskip{\normalsize \noindent Keywords: 
%% キーワード
Keywords in English. (approx. 5 Keywords)
}\vskip\baselineskip

%%%%%%%%%%%%%%%%%%%%%%%%%%%%%%%%%%%%%%%%%%%%%%%%%%%%%%%%
%% 送付日・著者所属
%%%%%%%%%%%%%%%%%%%%%%%%%%%%%%%%%%%%%%%%%%%%%%%%%%%%%%%%
\renewcommand{\thefootnote}{\fnsymbol{footnote}}
\footnotetext{\hspace{-4.5ex}
%% 送付日
Received on September 30, 2012.
}
\footnotetext{\hspace{-4.5ex}
%% 著者所属 1
$^{*}$ 産業技術大学院大学，産業技術研究科, School of Industrial Technology, AIIT
}
\footnotetext{\hspace{-4.5ex}
%% 著者所属 2
$^{**}$ 首都大学東京，システムデザイン研究科， 
Graduate School of System Design, Tokyo Metropolitan University
}

\renewcommand{\thefootnote}{*\arabic{footnote}}
\setcounter{footnote}{0}
\begin{multicols}{2} % 2段組の開始
\small % 本文を9ptにする

%%%%%%%%%%%%%%%%%%%%%%%%%%%%%%%%%%%%%%%%%%%%%%%%%%%%%%%%
%% ここから本文
%%%%%%%%%%%%%%%%%%%%%%%%%%%%%%%%%%%%%%%%%%%%%%%%%%%%%%%%
\section{はじめに}
本稿では産業技術大学院大学紀要のフォーマットについて記す．

\LaTeX のフォーマットを使う執筆者は
この``format.tex''ファイルの中身を書き換えて使うこと．
詳細な書式は``kiyou.sty''ファイルで定義しており，このファイルを変更してはならない．

\section{標題等ついて}
\subsection{原稿}
原稿は，日本語もしくは英語による完全版下(camera ready)原稿とする．製版後の校正は原則として不可能であるため，誤字や脱字がないよう，特に念を入れて仕上げる．刷り上がりは，6頁以上が望ましい．

\subsection{標題}
標題は和文ならびに英文とする．英文原稿の場合は，和文表題を記述するところに英文標題を記述し，通常の英文標題のところは削除すること．

\subsection{著者名・所属}
著者名も英文による原稿の場合は，通常の和文著者名のところに英文で記述し，英文著者名のところは削除すること．所属も英文で記述すること．

\subsection{アブストラクト・キーワード}
和文ではなく英文で記述すること．アブストラクトは100語程度とし，キーワードは5つ程度とする．

\subsection{標題等の割付}
別紙産技大紀要カメラレディ用見本に従って，[和文標題，和文著者名，英文標
題，英文著者名，英文アブストラクト，英文キーワード，所属]の割付を行う．

\section{本文について}

\subsection{余白}
天地左右余白(マージン)・段間余白(コラムスペース)も産技大紀要カメラレディ用見本に従う．

\subsection{見出し}
原稿には，大見出し，中見出しなどを設け，それらを明瞭に区分する．さらに細分を要するときは，著者の分類に委ねる．

\subsection{句読点}
句読点には，全角ピリオド(．)，全角コンマ(，)を用いること．


\section{図・表について}
\subsection{キャプション}
図・表には，図\ref{fig:AIIT}，図2，
表\ref{tab:cost2}，表2 のように論文全体で通し番号をつけること．
表のキャプションは表の上に，図のキャプションは図の下につけること．図・表
ともに配置は中央揃えにすること．英文の場合には，Fig. 1，Fig. 2，Table 1，
Table 2のように，番号をつけること．通し番号，標題は本文と同じ書体を使用
すること．

multicols を用いたため figure 環境，table 環境は使えない．
通常の caption が使えないため．代わりに \textbackslash figcaption，
\textbackslash tblcaption を用いて
図と表の caption をつけること．以下に例を示す．

\begin{center}
\includegraphics[width=0.3\textwidth]{AIIT_Symbol.eps}
\figcaption{図のキャプション}
\label{fig:AIIT}
\end{center}

\tblcaption{表のキャプション}
\label{tab:cost2}
\begin{center}
\begin{tabular}{|c|c|c|} \hline
A & B & C \\ \hline
$A_{1}$ & $B_{1}$ & $C_{1}$ \\ \hline
$A_{2}$ & $B_{2}$ & $C_{2}$ \\ \hline
$A_{3}$ & $B_{3}$ & $C_{3}$ \\ \hline
$A_{4}$ & $B_{4}$ & $C_{4}$ \\ \hline
\end{tabular}
\end{center}

\subsection{図の中の文字等}
 図・表は，印刷に十分耐えうるものでなければならない．刷り上がり時の文字が小さすぎないよう十二分に配慮し，線の太さにも注意する．

\subsection{色刷り}
図・表に色刷りを必要とする場合は，別途連絡すること．ただし，製本上の都合で色刷り頁を設けることができない場合もありうる．

\section{参考文献について}
参考文献は，通し番号とし，本文中では，当該事項または人名などの参考とする
後に，\cite{TRA96HuHaCo}，\cite{SICE02Yo}--\cite{Asa02Ar} のように記す．文章の末尾に記す必要がある場合には，句読点の前に記す．参考文献は，原則として，雑誌の場合は，著者，標題，雑誌名，巻，号，頁，年の順に記す．また，著書の場合は，著者，書名，発行所，発行年の順に記す．参考文献例を本文の最後に挙げるので参考されたい．

\section{おわりに}
本稿では産業技術大学院大学紀要のフォーマットについて記した．

\begin{thebibliography}{99}
\bibitem{TRA96HuHaCo} 
S. Hutchinson, G. D. Hager and P. I. Corke,
``A Tutorial on Visual Servo Control,''
{\it IEEE Trans. Robotics and Automation},
Vol.~12, No.~5, pp.~651--670, 1996.
%
\bibitem{SICE02Yo}
吉川恒夫,
``ロボット技術,''
計測と制御, Vol.~41, No.~1, pp.~17--21, 2002.
%
\bibitem{Joh06SpHuVi} 
M. W. Spong, S. Hutchinson and M. Vidyasagar, 
{\it Robot Modeling and Control}, 
John Wiley \& Sons, 2006. 
%
\bibitem{Asa02Ar}
有本 卓, 新版 ロボットの力学と制御, 朝倉書店, 2002. 
\end{thebibliography}

\newpage % 消さないで
\end{multicols}
\end{document}


